\documentclass[11pt]{article}
\usepackage{iftex}
\ifPDFTeX
  \usepackage[utf8]{inputenc}
  \usepackage[T1]{fontenc}
  % Portable ASCII renderings for symbols retained by the PDF text extractor.
  \DeclareUnicodeCharacter{00A7}{S}
  \DeclareUnicodeCharacter{00A8}{"}
  \DeclareUnicodeCharacter{00AF}{-}
  \DeclareUnicodeCharacter{00B4}{'}
  \DeclareUnicodeCharacter{00B5}{mu}
  \DeclareUnicodeCharacter{00B7}{.}
  \DeclareUnicodeCharacter{00D7}{x}
  \DeclareUnicodeCharacter{02C6}{hat}
  \DeclareUnicodeCharacter{0338}{/}
  \DeclareUnicodeCharacter{039B}{Lambda}
  \DeclareUnicodeCharacter{03B1}{alpha}
  \DeclareUnicodeCharacter{03B3}{gamma}
  \DeclareUnicodeCharacter{03B4}{delta}
  \DeclareUnicodeCharacter{03B5}{epsilon}
  \DeclareUnicodeCharacter{03B7}{eta}
  \DeclareUnicodeCharacter{03BA}{kappa}
  \DeclareUnicodeCharacter{03C0}{pi}
  \DeclareUnicodeCharacter{03C1}{rho}
  \DeclareUnicodeCharacter{03C6}{phi}
  \DeclareUnicodeCharacter{03D5}{phi}
  \DeclareUnicodeCharacter{2013}{-}
  \DeclareUnicodeCharacter{2014}{--}
  \DeclareUnicodeCharacter{2019}{'}
  \DeclareUnicodeCharacter{201C}{"}
  \DeclareUnicodeCharacter{201D}{"}
  \DeclareUnicodeCharacter{2022}{*}
  \DeclareUnicodeCharacter{2126}{Omega}
  \DeclareUnicodeCharacter{2190}{<-}
  \DeclareUnicodeCharacter{2192}{->}
  \DeclareUnicodeCharacter{2194}{<->}
  \DeclareUnicodeCharacter{21D2}{=>}
  \DeclareUnicodeCharacter{2206}{Delta}
  \DeclareUnicodeCharacter{2208}{in}
  \DeclareUnicodeCharacter{2212}{-}
  \DeclareUnicodeCharacter{2217}{*}
  \DeclareUnicodeCharacter{221A}{sqrt}
  \DeclareUnicodeCharacter{221D}{proportional-to}
  \DeclareUnicodeCharacter{2225}{||}
  \DeclareUnicodeCharacter{223C}{similar-to}
  \DeclareUnicodeCharacter{2248}{approximately}
  \DeclareUnicodeCharacter{2264}{<=}
  \DeclareUnicodeCharacter{2265}{>=}
  \DeclareUnicodeCharacter{226A}{<<}
  \DeclareUnicodeCharacter{226B}{>>}
  \DeclareUnicodeCharacter{2286}{subset-of}
  \DeclareUnicodeCharacter{22C6}{*}
\else
  \usepackage{fontspec}
  \setmonofont{Latin Modern Mono}[Scale=MatchLowercase]
\fi
\usepackage{geometry}
\usepackage{fancyvrb}
\geometry{margin=1in}
\RecustomVerbatimEnvironment{verbatim}{Verbatim}{
  fontsize=\scriptsize
}

\title{Beneventano--Poggio: Archival PDF Text Extraction}
\author{Converted from docs/theory/BP.pdf}
\date{PDF extracted mechanically for theorem audit}

\begin{document}
\maketitle

\noindent This file is a mechanical text extraction from \texttt{BP.pdf}. It is
kept in the paper archive for audit and comparison; it is not the original
LaTeX source.

\clearpage
\section*{Extracted page 1}
\begin{verbatim}
A Mathematical Framework for Building Agent Systems
A Mathematical Framework for Building Agent Systems
Task-First Design, Equivalence, and Model Choice under Deployment
Budgets
Pierfrancesco Beneventano and Tomaso Poggio
Center for Brains, Minds and Machines, MIT
Editor: Under review
Abstract
A builder of LLM agent systems faces a concrete design problem: given a task distribution
and a compute budget, which model, orchestration, and tool configuration minimizes time
to verified solution? We provide a mathematical framework that replaces benchmark racing—
the standard methodology of running both models and picking the winner—with structured
analysis.
The framework rests on three moves. First, we separate the task (an objective under
budgets with an external verifier) from the graph (the solver-induced execution), and show
this distinction changes design decisions. Second, we cast agent building as a stochastic
optimization problem over a six-dimensional design tuple (M, π, Ω, W, K, p) and prove that
restricting optimization to model variation alone is suboptimal whenever optimal routing
policies differ across models. Third, we define two notions of task equivalence—post-training
equivalence (when does training data transfer?) and orchestration equivalence (when does
the same workflow apply?)—and show they are independent.
For packed latent-mode families, expected certified log-speedup decomposes exactly as
∆ = log(κ /κ) + φ + G − ε, separating per-step cost, within-mode competence, task-relevant
0
information, and routing mismatch into four independently measurable terms. This yields a
closed-form crossover formula: a smaller specialist dominates a frontier generalist whenever
the effective mode count falls below a computable threshold. The decomposition is exact in
the packed-family regime and degrades gracefully outside it, with remainder O(δ log M +η).
n
We apply the framework to building agents for formal mathematical research, deriving
opposing prescriptions for different task regimes from the same theory, and provide a
measurement protocol for estimating all quantities from a structured pilot.
Keywords: agent design, certified time, task equivalence, orchestration equivalence, model
choice, operations research
1 Introduction
A builder of agent systems faces a question every deployment cycle: I have a task family,
a compute budget, and a menu of models ranging from fine-tuned 7B specialists to frontier
generalists. What system should I build? The question is not merely which model to call. It
is which model, which orchestration policy, which tools, which memory architecture, and
which budget allocation jointly minimize time to verified solution.
The dominant methodology is benchmark racing: run both models on a test set, pick
the winner, repeat when the next model arrives. Benchmark racing is fast and requires no
1
\end{verbatim}

\clearpage
\section*{Extracted page 2}
\begin{verbatim}
Beneventano and Poggio
theory. It is also structurally blind. It conflates per-step cost, within-mode competence,
task-relevant information, and routing quality into a single opaque comparison. A team that
benchmark-races cannot diagnose why model A beats model B, cannot predict when the
ranking reverses on a different task distribution, and cannot decide whether to fine-tune,
add tools, or restructure the workflow. Benchmark racing tells you the average winner. This
paper tells you the winner per task regime and tells you why.
Four questions. The paper is organized around four questions, ordered so that each
builds on the previous and the climax comes last.
1. What explains suboptimal performance? When an agent system is slow, is it
because the task lacks exploitable structure, because the system routes compute to the
wrong strategies, because each step is too expensive, or because the workflow graph is
badly shaped? The four-term decomposition (Section 6) separates these four causes into
independently measurable channels.
2. When are two tasks the same for design purposes? Benchmark racing treats
every new task as an independent evaluation problem. Post-training equivalence and
orchestration equivalence (Section 5) reveal when two apparently different tasks demand
the same system design—and when they do not.
3. When does training on one task help another? The transfer characterization
(Section 5) bounds the gain from source-task data by the mutual information between
the source representation and the target’s latent modes. Conditional guidance tells the
builder: shared modes implies transfer works; shared modes but different orchestration
implies retrain the router; no shared modes implies broader pre-training.
4. Which model should I use—and when does a small specialist beat a frontier
generalist? The crossover formula (Section 7) gives a closed-form threshold. A smaller
model wins when its cost advantage exceeds the information and routing penalties, with
the boundary determined by the task’s effective mode count.
Contributions. The paper makes five contributions.
1. A task-first ontology for agent systems (Section 3). We formalize budgeted trans-
ductive tasks, stochastic solvers, and certified time, separating the task (objective under
budgets) from the graph (solver-induced execution). A formal proposition shows this
separation changes design decisions: benchmark racing is suboptimal whenever optimal
routing differs across models.
2. Agent building as a deployment optimization problem (Section 4). We cast agent
design as stochastic optimization over a six-dimensional design tuple and prove that
restricting the program to model variation alone forfeits routing-adaptive designs.
3. Task equivalence and transfer (Section 5). We define post-training equivalence and
orchestration equivalence, prove they are independent, and characterize transfer gain.
√
The orchestration bound is O( α log M) via Pinsker and Audenaert.
4. An exact decomposition with approximation guarantee (Section 6). The four-
term decomposition is algebraically exact in the packed-family regime; outside it, a
2
\end{verbatim}

\clearpage
\section*{Extracted page 3}
\begin{verbatim}
A Mathematical Framework for Building Agent Systems
bounded remainder holds. The decomposition yields a crossover formula, a specialization-
dominance checklist, and a multi-agent threshold (Section 7).
5. A measurement protocol and application (Sections 8–9). We instantiate the
framework on formal mathematical research, deriving opposing prescriptions for different
task regimes, and provide estimators for all framework quantities from a structured pilot.
Scope. The main body concentrates on tasks for which success is externally checkable:
code against test suites, proofs against proof assistants, plans against constraint checkers.
This is a deliberate restriction, not an apology. It is precisely the regime in which time-to-
solution becomes mathematically clean and engineering-relevant. The exact decomposition
requires the packed latent-mode family (Assumption 6): each instance has a dominant
solution strategy, and certified time scales inversely with the compute share allocated to
that strategy. Outside this regime, the approximation guarantee (Proposition 25) quantifies
the degradation. Appendix A extends the framework to continuous rewards.
Relation to Achille–Soatto. Achille and Soatto [1] establish that the mutual informa-
tion between an agent’s history and its decisions bounds expected speedup over random
search—formally, that information buys speed. Their framework is model-independent: it
characterizes a task’s information structure without reference to which model processes it.
The moment one asks “which model should I use?”, one needs per-step cost and within-mode
competence—terms their framework does not contain. We extend their perspective: speed
has a price, the price depends on design choices, and the price structure is decomposable
into four measurable channels. They diagnose how much structure a task contains. We tell
you which system to build to exploit that structure.
Roadmap. Section 2 positions the paper. Section 3 defines the task, solver, and certified-
time objects. Section 4 formulates the deployment optimization problem. Section 5 develops
task equivalence and transfer. Section 6 contains the decomposition core. Sections 7–9 derive
consequences, apply the framework to mathematical research, and specify the measurement
protocol. Section 10 discusses limitations and open problems.
2 Related Work and Positioning
Achille–Soatto and the time-centric transductive view. The conceptual starting
point is Achille–Soatto [1], who frame agents in the verifiable transductive regime and
make time-to-solution central. We inherit that move almost entirely. Our contribution
is narrower and more operational: we specialize the viewpoint to the problem of building
agent systems, with explicit objects for tasks, solvers, budgets, transfer, orchestration
mismatch, and deployment-facing model choice. The older lineage behind this viewpoint
runs through Solomonoff induction, Levin search, and related resource-bounded views of
inference [37, 24, 29, 25, 34]. These works motivate the idea that prior weight, description
length, and reusable structure can be traded against search time. The present paper stays
within that lineage but asks a deployment-facing question: given a modern agent stack with
LLMs, tools, memory, and orchestration, which quantities should a builder optimize?
Graph-centric reasoning and workflow literature. A substantial body of recent work
describes reasoning or orchestration through explicit graph-like objects: trees of thoughts,
3
\end{verbatim}

\clearpage
\section*{Extracted page 4}
\begin{verbatim}
Beneventano and Poggio
graphs of thoughts, and code-represented workflows searched over graph-structured execution
spaces [44, 6, 45]. We agree that graphs are often the right representation of a solver execution.
Our difference is ontological rather than adversarial: the graph is not the task. The task
is the objective under budgets; the graph is the realization induced by a particular solver,
model, and memory regime. That separation is what lets us ask whether different solvers
on the same task should induce different graphs, and whether apparently similar graphs
correspond to the same task at all.
Algorithm portfolios and anytime computation. The question “how should one
allocate limited compute across competing strategies for a given instance?” has a long
history in algorithm portfolios and per-instance selection [20, 43], anytime algorithms and
metareasoning under deadlines [47, 19, 33], and sequential inspection [42, 40, 10]. Weitzman’s
Pandora’s Box result gives optimal stopping rules for inspecting alternatives with known
reward distributions—structurally close to our routing problem, but in a static setting
without the learning-from-verification feedback that the packed-family model captures. The
present framework brings this family of ideas into the language of modern agent systems:
the builder allocates compute across branches, tools, and strategies, and the target is a
certified time frontier rather than a scalar accuracy metric.
Operations research and team decision problems. Radner [31] initiated the theory
of team decisions: agents with a common objective but decentralized information. Our
multi-agent threshold (Section 7) is a special case of Marschak–Radner team optimization [28]
applied to the packed-family setting. The broader framing of agent building as a stochastic
optimization over design variables follows the tradition of stochastic programming [7] and
automated ML pipelines [21, 16]. The contribution is not the optimization framing—that is
standard—but the decomposable objective with four named, measurable terms.
Conditional computation and routing. Mixture-of-experts and conditional computa-
tion architectures separate a predictor from a router and make routing quality a first-class
determinant of performance [22, 35, 14]. This is structurally close to our mismatch account-
ing. The difference is granularity: in this paper the objects being routed are not token-level
experts but agentic modes—toolchains, decomposition choices, branches of search, and
orchestration policies.
Task equivalence and MDP similarity. Ferns, Panangaden, and Precup [15] define
bisimulation metrics for MDPs; Ravindran and Barto [32] introduce MDP homomorphisms.
Both operate at a finer grain than our equivalence notions: they compare state-action spaces,
while we compare posterior mode distributions at the task-family level. Agarwal et al. [2]
propose policy similarity metrics. Taylor and Stone [39] survey transfer learning in RL. Our
post-training equivalence addresses a different question: not whether two MDPs admit the
same optimal policy, but whether training data from one task family improves certified time
on another. Orchestration equivalence is, to our knowledge, without close precedent at this
grain size.
Verifiable reasoning and test-time compute. The motivating workloads are those in
which test-time reasoning is evaluated against an external checker: code generation against
unit tests, formal theorem proving, and verifier-backed reasoning [9, 27, 46]. A parallel
4
\end{verbatim}

\clearpage
\section*{Extracted page 5}
\begin{verbatim}
A Mathematical Framework for Building Agent Systems
literature studies how additional test-time compute improves performance [41, 44, 36]. Our
contribution is not another prompting or search heuristic. It is a mathematical language for
asking what kind of speedup a design has bought and whether that speedup came from more
reusable structure, better orchestration, or simply more expensive per-step computation.
Memory, retrieval, and mode discovery. Retrieval-augmented models and library-
induction systems convert past experience into faster future execution [26, 8, 13, 18]. In
our language, such mechanisms change both how much reusable structure is available and
how effectively the solver deploys it. Mode discovery—identifying the latent strategies that
govern performance—connects to the options framework [38], option-critic architectures [4],
and trajectory-based behavioral mode discovery [5]. Our measurement protocol (Section 9)
adapts trajectory clustering to the proof-strategy setting.
Positioning. We do not claim priority on time, verifiability, search, routing, or graph-based
workflows in isolation. The claim is that, within one framework, we separate tasks from
solver-induced graphs, cast agent building as a budgeted design problem, formalize task
equivalence and transfer, and connect these objects to theorem-level speed accounting and
deployment-facing model choice. The paper is a framework for building agent systems, not
only for interpreting them after the fact.
3 A Mathematical Framework for Agentic Tasks
The mathematical objects in this paper are tasks, solvers, and certified time. The definitions
are compressed; a Lean 4 tactic-prediction example follows each one to build intuition.
3.1 Tasks
Definition 1 (Budgeted transductive task) A budgeted transductive task at scale n
is a tuple T = (X , Y , V , B ) where:
n n n n
• X is a set of instances and Y a set of candidate solutions;
n n
• V : X × Y → {0, 1} is a polynomial-time verifier;
n n n
• B = (b , b , b ) specifies budgets for time, cost, and memory.
n t c m
Success on instance x means producing y with V (x, y) = 1 within B .
n n
In Lean 4: X is a set of lemma statements, Y is the set of tactic proof scripts, V is
n n n
the Lean kernel checker, and B specifies maximum API cost and wall-clock time.
n
Definition 2 (Task family) A task family is a sequence T = (T ) of budgeted trans-
n n≥1
ductive tasks, together with a deployment distribution µ over instances X . The family
n n
captures a class of problems the builder expects to encounter repeatedly.
The Lean example: a task family is a distribution over Mathlib lemma difficulties—some
routine, some library-heavy, some creative.
5
\end{verbatim}

\clearpage
\section*{Extracted page 6}
\begin{verbatim}
Beneventano and Poggio
3.2 Solvers and orchestration graphs
Definition 3 (Stochastic solver) A stochastic solver A for task family T is a randomized
algorithm that, given instance x ∈ X and budgets B , produces a candidate y (or reports
n n
failure) using a sequence of model calls, tool invocations, memory accesses, and branching
decisions. A run of A on (x, r) (where r denotes randomness) induces a run graph G (x; r)—
A
a DAG whose nodes are computational steps and whose edges encode data and control
flow.
The graph is not the task. The task is the objective under budgets; the graph is the
realized execution induced by a particular solver, model, and orchestration on a particular
instance. Two solvers for the same task may produce radically different graphs—one a linear
chain of rewrites, another a branching tree search with backtracking—yet both target the
same objective. Confusing the graph with the task is the structural error that benchmark
racing commits: it evaluates the realized execution without decomposing the objective.
3.3 Certified time
Definition 4 (Certified time-to-solution) For solver A on task family T at confidence
level δ ∈ (0, 1), the certified time-to-solution is
T (n, δ) = inf{t > 0 : P [A(x; r) succeeds within time t] ≥ 1 − δ} .
A x∼µn,r
Certified time is the deployment-relevant quantity: the builder needs the system to
succeed with probability 1 − δ within a time budget, not merely to succeed “on average.”
The log of certified time, log T , is the natural scale for additive decomposition.
A
Definition 5 (Certified log-speedup) Given a baseline solver A and a deployed solver
0
A, the certified log-speedup is ∆ = log T (n, δ) − log T (n, δ). Positive ∆ means the
A0 A
deployed solver is faster.
3.4 The packed latent-mode family
The exact decomposition requires a structural assumption on the task family. The assumption
is stylized but captures the essential feature of verifiable tasks with multiple solution
strategies.
Assumption 6 (Packed latent-mode family) Fix a task family at scale n. Assume:
1. Latent modes. There is a discrete latent variable S ∈ [M ] := {1, . . . , M } with prior π,
n n
representing the solution strategy, and an observable D ∼ P (D | S) representing verifier
feedback and partial execution traces, with posterior π (s) := P[S = s | D].
D
2. One expert per mode. Each mode s has a specialized expert with certified internal scale
t (s).
0
3. Allocation determines time. A routing policy chooses a share vector q ∈ ∆ . Under
Mn
hidden mode s, allocating share q(s) yields certified time
t (s)
0
T (s) = c · ,
q sim
q(s)
6
\end{verbatim}

\clearpage
\section*{Extracted page 7}
\begin{verbatim}
A Mathematical Framework for Building Agent Systems
where c is the simulation overhead.
sim
4. Baseline and deployed schedulers. The baseline chooses a fixed share vector q ∈ ∆ ;
0 Mn
the deployed system chooses a data-dependent map D (cid:55)→ q ∈ ∆ .
D Mn
In Lean tactic prediction, the modes are proof strategy families: induction, case splitting,
library search, simp chains, rewrite sequences. The observable D is the partial execution
trace—which tactics were tried, which failed, what error messages the kernel returned. The
allocation q(s) is the fraction of the attempt budget devoted to strategy s. The assumption
that T (s) ∝ t (s)/q(s) is a geometric-trials model: doubling the share on the correct
q 0
strategy halves expected time to find the proof. This is well-supported empirically for
search-based theorem proving [17].
The entropy-effective mode count. For a posterior π over modes, define M (D) :=
D eff
exp(H(π )). When M is small on most instances, only a few modes need to be separated
D eff
and executed well—the regime where specialization is most valuable.
3.5 The graph is not the task
The following proposition formalizes why the task-graph separation matters for design
decisions.
Proposition 7 (Benchmark racing can be suboptimal) Under the packed-family as-
sumption, consider two models M and M with distinct optimal routing policies: π∗,A ̸=
A B D
π∗,B on a set of instances of positive measure. Then evaluating both models under a single
D
fixed workflow (as benchmark racing does) can select the wrong model: the model with lower
win rate under the fixed workflow may achieve lower certified time under its own optimal
routing.
Proof Model M evaluated under its optimal routing has mismatch ε (π∗,A) = 0. Model
A A
M evaluated under M ’s routing has mismatch ε (π∗,A) = E [KL(π∗,B∥π∗,A)] > 0.
B A B D D D
Benchmark racing compares both under π∗,A, penalizing M by ε (π∗,A) nats. Under
B B
its own routing, M recovers the full penalty. If M ’s other advantages exceed M ’s by
B B A
less than ε (π∗,A) but more than zero, benchmark racing selects A while B with adapted
B
routing is superior.
A concrete instance: Model A wins 60% of problems with a fixed workflow; Model B
wins only 45% with the same workflow but 70% with an adapted routing policy. Benchmark
racing selects A; the framework selects B. The graph is not the task. The task is the same
in both cases; the graphs—and therefore the performance—differ because the routing policy
differs.
Scope declaration. The framework restricts to verifiable tasks with external checkers.
This already includes many of the most important agentic workloads: code synthesis
against tests, formal theorem proving, constrained planning, and verifier-backed search.
The restriction is a feature: it is precisely the regime in which time-to-solution becomes
mathematically clean and engineering-relevant.
7
\end{verbatim}

\clearpage
\section*{Extracted page 8}
\begin{verbatim}
Beneventano and Poggio
4 Agent Building as a Deployment Optimization Problem
The task object and solver-induced graph are now in place. The question is no longer only
why a finished system is slow. It is how to choose model, memory, tools, routing, and
parallelism before deployment.
4.1 The design tuple
A design point is
d = (M, π, Ω, W, K, p) ∈ D , (1)
n
where M is the base model, π is the orchestration/routing policy, Ω is the available tool
library, W denotes long-term memory resources, K is the active context budget, and p is
the allowed degree of parallelism. Write A for the solver induced by design d.
d
In the Lean example: M is DeepSeek-Prover-7B or Claude Opus; π is the tactic-selection
and backtracking policy; Ω includes the Lean kernel, Mathlib search, and rewriting tools; W
is the persistent lemma cache; K is the proof-state context window; and p is the number of
parallel proof branches.
4.2 The master optimization problem
For a confidence level δ ∈ (0, 1) and a deployment loss Λ over feasible time–cost–memory–
n
parallelism tuples, the builder solves
min Λ (t, c, m, p) subject to (t, c, m, p) ∈ Fdep(n, δ) (2)
n A
d∈Dn d
where Fdep(n, δ) is the deployment frontier: the set of time–cost–memory–parallelism tuples
A
d
at which solver A succeeds with probability ≥ 1 − δ.
d
This is the sense in which building an agent system is an operations-research problem.
The builder chooses a design point under resource constraints, not merely a single model
under a scalar accuracy metric. The paper does not claim to solve this optimization in full
generality. It identifies the design-variable space, provides a decomposable objective for the
model-choice component, and solves restricted sub-problems (Sections 6–7).
4.3 Suboptimality of model-restricted optimization
Benchmark racing restricts the design program to model variation alone: it fixes the routing
policy and all other coordinates, then compares models. The following proposition quantifies
what this restriction forfeits.
Proposition 8 (Suboptimality of M-restricted optimization) Under the packed-
family assumption (Assumption 6), consider two candidate models M and M with distinct
A B
optimal routing policies: π∗,A ≠ π∗,B on a set of instances of positive measure. If the builder
D D
restricts to M-variation alone—fixing π = π∗,A and all other coordinates at model A’s
values—then the suboptimality gap relative to the unrestricted optimum is at least
ε (π∗,A) − ε (π∗,B) = ε (π∗,A) > 0,
B B B
the routing-mismatch penalty from evaluating model B under the wrong routing policy.
8
\end{verbatim}

\clearpage
\section*{Extracted page 9}
\begin{verbatim}
A Mathematical Framework for Building Agent Systems
Proof Under the four-term decomposition (Theorem 21), the certified log-speedup
for model M under routing policy q is ∆(M , q) = log(κ /κ ) + φ + G − ε (q),
B B 0 B B B B
where ε (q) = E [KL(π∗,B∥q )]. The unrestricted optimum uses q = π∗,B, achieving
B D D D D
ε (π∗,B) = 0. Under M-restriction, the builder evaluates M with π∗,A, incurring
B B
ε (π∗,A) = E [KL(π∗,B∥π∗,A)] > 0 since the policies differ on positive measure and KL
B D D D
divergence is zero if and only if the distributions agree almost surely.
Corollary 9 (Benchmark racing eliminates routing-adaptive designs) Benchmark
racing can only discover which model is better under one fixed routing policy. It cannot
discover that model B with an adapted routing policy outperforms model A—precisely the
design the framework identifies as optimal in the contested zone where M ≈ M⋆ .
eff eff
Builder interpretation. The suboptimality gap is largest when two models’ optimal
routing policies differ substantially—the contested zone where the framework’s marginal value
over benchmark racing is concentrated. In the clear-cut zones (M ≪ M⋆ or M ≫ M⋆ ),
eff eff eff eff
one model dominates regardless of routing, and benchmark racing gives the right answer.
4.4 What does posing the problem this way buy?
Without the OR framing, the builder does trial-and-error over a combinatorial design space.
With it, the builder has structured search with named dimensions: model choice affects
cost and competence; routing policy affects mismatch; tools and memory affect available
structure; context budget trades node cost against graph shrinkage; parallelism trades
coordination overhead against critical-path reduction. The four-term decomposition provides
the objective function for the model-choice component.
Which sections address which coordinates. Model choice M is addressed by Section 7
(crossover theorem under simplifying assumptions). Routing policy π is partially addressed
by Section 5 (orchestration equivalence: when two tasks demand the same routing) and
Section 6 (the ε term diagnoses routing mismatch). Transfer across task distributions is
characterized by Section 5. Orchestration architecture Ω, context budget W , and parallelism
p are addressed only partially. Full joint optimization over all six coordinates simultaneously,
multi-agent team composition beyond the threshold result, and optimal budget allocation
under cost constraints remain open problems.
Closing promise. The design program identifies six decision variables. The next two
sections show that three of these variables are often redundant—because tasks that appear
different are, for design purposes, the same—and that the remaining objective decomposes
exactly into four terms a builder can estimate and independently improve.
5 When Are Two Tasks the Same?
Benchmark racing treats every new task as an independent evaluation problem. The
equivalence framework reveals when two apparently different tasks are the same design
problem—and when training on one transfers to the other.
9
\end{verbatim}

\clearpage
\section*{Extracted page 10}
\begin{verbatim}
Beneventano and Poggio
5.1 Budget-preserving reduction and task equivalence
Definition 10 (Budget-preserving reduction) A budget-preserving reduction from
task family TA to task family TB at scale n is a pair of maps (f, g) where f : X A → X B
n n
maps instances and g : YB → YA maps solutions, such that:
n n
1. Correctness preservation: V B(f(x), y) = 1 =⇒ V A(x, g(y)) = 1.
n n
2. Budget preservation: the composition x (cid:55)→ g(AB(f(x))) respects the budget constraints
d
of TA: the overhead of applying f and g is absorbed within BA.
n
Definition 11 (Strong and weak equivalence) Two task families TA and TB are:
• Strongly equivalent if budget-preserving reductions exist in both directions with overhead
that is o(1) relative to the budgets.
• α-weakly equivalent if for every design d the certified log-speedups satisfy |∆A(d) −
∆B(d)| ≤ α.
Strong equivalence says the two tasks are the same design problem: any solver for one is
a solver for the other at negligible overhead. Weak equivalence is looser: the two tasks may
differ, but no design choice exploits the difference by more than α nats.
5.2 Post-training equivalence
The builder’s question: if a model is post-trained on task family A and deployed on family B,
should the builder expect the same performance?
Definition 12 (Post-training efficiency) Fix a base solver U(0) and an update rule Upd.
For data budget m, the post-training efficiency of A-data for task B is the transfer ratio
∆ (n, m, δ)
A→B
η (m) := .
A→B
∆ (n, m, δ)
B→B
When η ≈ 1, training on A-data is almost as good as training on B-data for improving
A→B
performance on B.
Theorem 13 (Post-training equivalence from task equivalence) In the packed-
family regime, suppose both families share the same latent space S ∈ [M ] with different
n
observation channels P (D|S) and P (D|S). Then
A B
(m)
I(S; D ) − ε
η (m) = A A→B . (3)
A→B
(m)
I(S; D ) − ε
B B→B
In particular:
(i) If both families share all latent structure and cross-family mismatch equals within-family
mismatch, then η = 1: training on A is exactly as good as training on B.
A→B
(ii) If the families share the latent space but different channels, then η ≈ I(S; D )/I(S; D )
A→B A B
when mismatch terms are small: the transfer ratio is governed by the information-quality
ratio of the two channels.
10
\end{verbatim}

\clearpage
\section*{Extracted page 11}
\begin{verbatim}
A Mathematical Framework for Building Agent Systems
Proof Apply the Routing-Information Identity (Theorem 19) to both transfer settings over
(m)
the shared mode space. For the numerator, the side-information variable is D drawn
A
from the A-channel, applied to target family B:
∆ = I(S; D (m) ) − ε , where ε := E (cid:2) KL (cid:0) πB ∥q (· | D ) (cid:1)(cid:3) .
A→B A A→B A→B DA DA B A
Here πB (s) = P[S = s | D ] is the posterior over B-task modes given A-channel data. For
DA A
(m)
the denominator, within-family data gives ∆ = I(S; D ) − ε . Dividing yields (3).
B→B B B→B
Claims (i) and (ii) follow by specialization.
The shared-latent-space assumption is strong. It is appropriate when both families
draw on the same underlying skill bank—coding and formal mathematics both involve
decomposition, symbolic manipulation, and verification. It is inappropriate when the
families share only surface-level features.
Design heuristic for mixed-source post-training. For fraction α of data from source A,
(m)
the natural objective is ∆ (m) = I(S; D ) − ε . The mutual information is concave
α→B α α
in α by standard arguments [11]. The behavior of ε is less clear: an earlier draft asserted
α
convexity by citing “the infimum of convex functions is convex,” which is false in general.
We state as a design heuristic: when the concavity of I(S; D ) dominates any non-convexity
α
in ε , compute ∆ on a coarse grid and select the maximizer. Whether the objective is
α α→B
generically concave under natural structural conditions is an open question (Section 10).
5.3 Transfer characterization
Proposition 14 (Transfer through target-relevant side information) Suppose the
target family TB admits the packed-family representation with latent mode S ∼ πB. Let
B
adaptation on family TA produce a random representation Z , and let the transferred system
A
choose allocation q (· | Z ). Write πB (s) := P[S = s | Z ]. Then the packed-family
B A ZA B A
certified-time gain on TB from using Z is
A
∆packed = I (S : Z ) − E(cid:2) KL (cid:0) πB ∥q (· | Z ) (cid:1)(cid:3) . (4)
A→B πB B A ZA B A
Proof Apply the packed-family decomposition to the target family TB with side-information
variable Z rather than within-family data. The proof is otherwise unchanged: the cross-
A
entropy identity splits the allocation term into information and mismatch.
Remark 15 (DPI bound under joint latent structure) If additionally D ← S ↔
A A
S → D is a Markov chain and Z is a function of (D , S ), then by the data-processing
B B A A A
inequality: ∆packed ≤ I(S : S ). Transfer is bounded by the mutual information between the
A→B A B
latent mode spaces. The Markov chain assumption is appropriate when the families share a
common underlying skill space; it is inappropriate when they share only surface features.
11
\end{verbatim}

\clearpage
\section*{Extracted page 12}
\begin{verbatim}
Beneventano and Poggio
Estimability. The quantity I(S ; Z ) is a theoretical characterization, not an actionable
B A
prescription: estimation requires model internals or representation probing. For practical
guidance, the builder uses the three-case conditional heuristic below and an estimable
mode-overlap proxy—the Bhattacharyya coefficient between empirical mode distributions
from two pilot runs. Estimation of I(S ; Z ) is an open problem (Section 10).
B A
Conditional guidance. The framework yields three conditional prescriptions for transfer:
• Shared dominant modes, similar orchestration. Training benefit transfers. Deploy the
source-trained model with minor adaptation.
• Shared dominant modes, different orchestration. Retrain the routing policy but keep the
base model. Competence transfers; the routing does not.
• Target has modes not present in the source. Broader pre-training matters more than
source-task fine-tuning. The source representation carries little information about the
target’s latent modes.
5.4 Orchestration equivalence
Post-training equivalence governs what model to train. Orchestration equivalence governs a
complementary question: when to reuse the same workflow.
Definition 16 (α-Orchestration equivalence) Two packed-family task families T A and
T B sharing a common mode bank [M ] are α-orchestration-equivalent if
n
E D (cid:2) KL (cid:0) π D A (cid:13) (cid:13) π D B(cid:1)(cid:3) ≤ α, (5)
where πA, πB ∈ ∆ are the posterior mode distributions induced by side information D
D D Mn−1
under the two families, and the expectation is over the marginal distribution of D.
Two families that are α-orchestration-equivalent with small α demand similar routing
policies: the posterior beliefs about which strategy to deploy are close on average.
Proposition 17 (Orchestration equivalence bounds) Suppose T A and T B are α-
orchestration-equivalent packed families with common mode bank [M ]. Then:
n
(a) Information-term bound (independent of routing policy q).
|G A − G B | = (cid:12) (cid:12) E D [H(π D A)] − E D [H(π D B)] (cid:12) (cid:12) ≤ √ 2α log (cid:18) √ M n (cid:19) , (6)
2α
√
provided 2α ≤ 1 − 1/M .
n
(b) Routing-mismatch bound (depends on routing policy q). For any routing policy q
with q := min q(s) > 0:
min s
√ √ (cid:18) M (cid:19)
n
|ε (q) − ε (q)| ≤ 2α log(1/q ) + 2α log √ . (7)
A B min
2α
(c) Simplified bound under uniform-floor routing. If q ≥ 1/M , both bounds are
√ min n
O( α log M ).
n
12
\end{verbatim}

\clearpage
\section*{Extracted page 13}
\begin{verbatim}
A Mathematical Framework for Building Agent Systems
Proof sketch. Apply Pinsker’s inequality [30, 12] pointwise: ∥πA − πB∥ ≤
D D 1
(cid:113) √
2 KL(πA∥πB). Take expectations and apply Jensen (concavity of ·): δ¯ :=
D D √
E[∥πA − πB∥ ] ≤ 2α. For (a), apply Audenaert’s [3] entropy continuity lemma:
D D 1
|H(p) − H(r)| ≤ ∥p − r∥ log(M /∥p − r∥ ) for distributions on M atoms. Take expecta-
1 n 1 n
tions; t (cid:55)→ t log(M /t) is concave, so Jensen gives the bound in terms of δ¯. For (b), use the
n
KL-difference decomposition:
(cid:80)
KL(p∥q) − KL(r∥q) = [H(r) − H(p)] + [p(s) − r(s)] log(1/q(s)).
s
Bound the first term by (a) and the second by δ¯ · log(1/q ) via Pinsker. Full proof in
min
Appendix B.
Builder interpretation. If two tasks are α-orchestration-equivalent, the penalty from
√
reusing task A’s routing policy on task B is at most O( α log M ) nats in both the
n
information and routing terms. For a system with M = 10 modes and α = 0.01, this
√ n
is about 2 0.02 · log(10) ≈ 0.65 nats—less than the typical cost-ratio term log(30) ≈ 3.4
√
nats. The bound degrades with α, not α, reflecting the sub-linear relationship between
KL divergence and total-variation distance. This is weaker than a linear bound for small α
but is the correct scaling from first principles.
5.5 The two equivalences are independent
The two equivalence notions capture independent design dimensions.
Orchestration-equivalent but not post-training-equivalent. Two families can share
the same posterior structure (same beliefs about which mode to use, hence small α) while
one requires much more data to reach those beliefs. The posteriors are close; the data
requirements are not. A builder can reuse the workflow but must invest in separate training.
Post-training-equivalent but not orchestration-equivalent. Two families can share
the same latent modes (training on one transfers to the other, so η ≈ 1) while the
A→B
posterior structures differ enough that they demand different routing policies. Coding and
formal mathematics may share symbolic-manipulation modes, so fine-tuning transfers. But
the optimal proof-search orchestration differs from the optimal code-repair orchestration. A
builder can reuse the model but must redesign the workflow.
This separation is not a pathology. It is the conceptual payoff of having two distinct
equivalence notions: post-training equivalence tells the builder when to reuse training data;
orchestration equivalence tells the builder when to reuse the workflow. The decomposition
in the next section provides the objective these equivalences operate on.
6 Information-Routing and Model-Aware Decompositions
The equivalence notions of the previous section identify when two tasks are the same
design problem. The present section provides the objective function that design problem
optimizes. We first establish the Routing-Information Identity—the two-term decomposition
∆ = G−ε that is the paper’s cleanest equation—then extend it to the four-term Model-Aware
Decomposition that compares models, and finally prove the upper envelope theorem and the
approximation guarantee outside the packed family.
13
\end{verbatim}

\clearpage
\section*{Extracted page 14}
\begin{verbatim}
Beneventano and Poggio
6.1 The Routing-Information Identity
The first result decomposes certified log-speedup into two terms: how much the system
learns about the latent mode, and how well it uses what it learns.
Lemma 18 (Conditional log-time decomposition) Under Assumption 6, for any re-
alized feedback D and any share vector q ∈ ∆ with supp(π ) ⊆ supp(q),
Mn D
E (cid:2) log T (S) | D (cid:3) = E [log(c t (S))] + H(π ) + KL(π ∥q).
S∼πD q S∼πD sim 0 D D
The allocation q = π uniquely minimizes conditional expected certified log-time. The excess
D
over this optimum is exactly KL(π ∥q).
D
Proof Since log T (s) = log(c t (s)) − log q(s),
q sim 0
E (cid:2) log T (S) | D (cid:3) = E [log(c t (S))]+ (cid:88) π (s) log 1 = E [log(c t (S))]+H(π )+KL(π ∥q),
S∼πD q S∼πD sim 0 D
q(s)
S∼πD sim 0 D D
s
(cid:80)
using the cross-entropy identity p(s) log(1/q(s)) = H(p) + KL(p∥q). Nonnegativity of
s
KL divergence gives the optimality claim.
Theorem 19 (Routing-Information Identity) Under Assumption 6, with a prior-
matched baseline q = π,
0
∆ = G − ε
where
G := I (S : D) = H(π) − E [H(π )], (8)
π D D
ε := E [KL(π ∥q )] . (9)
D D D
The term G is the mutual information between the latent mode and the feedback—the
information generation quality. The term ε is the expected KL divergence of the routing
policy from the posterior—the routing mismatch. Equivalently: ∆ ≤ G, with equality if and
only if q = π almost surely. One nat of routing mismatch costs exactly one nat of certified
D D
log-speedup.
Proof Apply Lemma 18 to the baseline allocation q = π:
0
E(cid:2) log T (S) (cid:3) = E [log(c t (S))] + H(π).
π π sim 0
Apply it to the deployed allocation q :
D
E(cid:2) log T (S) (cid:3) = E [log(c t (S))] + E [H(π )] + E [KL(π ∥q )].
qD π sim 0 D D D D D
Subtract:
∆ = H(π) − E [H(π )] − E [KL(π ∥q )] = G − ε.
D D D D D
14
\end{verbatim}

\clearpage
\section*{Extracted page 15}
\begin{verbatim}
A Mathematical Framework for Building Agent Systems
The identity is algebraic: no approximation, no limit, no asymptotic qualifier. In the
packed-family regime, information generation and routing mismatch are the only deter-
minants of speedup, and they combine additively on the log scale. A system can be fast
because it generates good information (G large) or slow because it routes poorly (ε large),
and these two failures are separable.
For a non-prior-matched baseline q ̸= π, the identity becomes ∆ = G + KL(π∥q ) − ε,
0 0
adding a baseline suboptimality bonus. The prior-matched case is the clean comparison.
6.2 The Model-Aware Decomposition
The Routing-Information Identity holds model fixed. To compare models, we extend the
packed-family assumption to include model-dependent costs and competence.
Assumption 20 (Model-dependent packed family) Fix a task family and a shared
latent mode variable S ∼ π over [M ]. For each model M:
n
1. The per-step cost is κ(M).
2. The within-mode certified scale is t (M, s), so that certified time under mode s with
0
allocation share q(s) is
t (M, s)
(M) 0
T (s) = κ(M) · c · .
q sim q(s)
3. The feedback observable is D ∼ P (D | S), with posterior πM(s) := P[S = s | D = D].
M D M
4. The routing policy is qM ∈ ∆ .
D Mn
Theorem 21 (Model-Aware Decomposition) Under Assumption 20, comparing a
prior-matched baseline using model M against a deployed system using model M,
0
κ(M )
0
∆ = log + φ(M , M) + G(M) − ε(M) (10)
M 0
κ(M)
(cid:124) (cid:123)(cid:122) (cid:125) (cid:124) (cid:123)(cid:122) (cid:125) (cid:124) (cid:123)(cid:122) (cid:125)
(cid:124) (cid:123)(cid:122) (cid:125) competence information routing mismatch
per-step cost
where the four terms are:
κ(M )
0
log (per-step cost advantage of M), (11)
κ(M)
(cid:20) (cid:21)
t (M , S)
φ(M , M) := E log 0 0 (within-mode competence advantage), (12)
0 π
t (M, S)
0
G(M) := I (S : D ) (information generation quality), (13)
π M
ε(M) := E (cid:2) KL(πM∥qM) (cid:3) (routing mismatch). (14)
D D D
Proof Write out the log-time:
(M)
log T (s) = log κ(M) + log c + log t (M, s) − log q(s).
q sim 0
15
\end{verbatim}

\clearpage
\section*{Extracted page 16}
\begin{verbatim}
Beneventano and Poggio
For the prior-matched baseline with model M :
0
(cid:104) (cid:105)
E log T (M0) (S) = log κ(M ) + log c + E [log t (M , S)] + H(π).
π 0 sim π 0 0
For the deployed system with model M:
(cid:104) (cid:105)
E log T (M) (S) = log κ(M) + log c + E [log t (M, S)]
qD sim π 0
+ E [H(πM)] + E [KL(πM∥qM)],
D D D D D
where the last two terms arise from the cross-entropy identity exactly as in Theorem 19.
Subtracting and grouping:
(cid:20) (cid:21)
∆ = log κ(M 0 ) + E log t 0 (M 0 , S) + (cid:0) H(π) − E [H(πM)] (cid:1) − E [KL(πM∥qM)]
M κ(M) π t (M, S) D D D D D
0
κ(M )
0
= log + φ(M , M) + G(M) − ε(M).
0
κ(M)
The four terms have clean interpretations:
• Per-step cost log(κ /κ): the log cost ratio. A 7B model at 1/30th the price contributes
0
log 30 ≈ 3.4 nats.
• Competence φ: the expected log-ratio of within-mode solve times. A fine-tuned specialist
that is twice as fast per mode contributes log 2 ≈ 0.7 nats.
• Information G: the mutual information between the latent mode and verifier feedback.
A task with 8 equiprobable modes and a verifier that identifies the correct one contributes
log 8 ≈ 2.1 nats.
• Routing mismatch ε: the KL divergence of the routing policy from the posterior. A
policy that spreads mass uniformly when the posterior is concentrated wastes ε nats of
potential speedup.
Benchmark racing reports only ∆. The decomposition reports why ∆ takes the value it
does. Two systems with the same ∆ may differ entirely in how they achieve it—one through
low cost, the other through good routing—and the appropriate intervention is different in
each case.
6.3 The upper envelope theorem
The decomposition is exact within the packed family. Before specializing further, we state
the general upper bound that holds outside it. The theorem rests on the following regularity
conditions on the solver space and a structural witness.
Assumption 22 (A1–A7: shared admissibility and aligned witness)
• A1 (verifiability): V ∈ P.
n
16
\end{verbatim}

\clearpage
\section*{Extracted page 17}
\begin{verbatim}
A Mathematical Framework for Building Agent Systems
• A2 (prefix coding): A is countable and prefix-free, and both w(·) and w(· | D ) satisfy
n n
Kraft.
• A3 (shared admissibility): the same nonempty set Aδ is used in the baseline and condi-
n
tioned schedule objectives.
• A4 (shared simulation): both schedules incur the same simulation overhead c (n) ∈
sim
poly(n).
• A5 (resource-bounded conditioning): the map D (cid:55)→ w(· | D ) is computable or decodable
n n
in polynomial time.
• A6 (structure witness class): there is a nonempty subclass S(H ) ⊆ Aδ of admissible
n n
solvers that realize the reusable structure H .
n
• A7 (aligned witness): there exists a⋆ ∈ S(H ) and nonnegative slacks (ρ , ρ , κ ) such
n U D U
that
F (a⋆) ≤ eρU inf F (a), F (a⋆) ≤ eρD inf F (a),
U U D D
a∈Aδ a∈Aδ
n n
− log w(a⋆) ≤ Kpoly(H ) + κ ,
n U
where F and F denote the baseline and conditioned schedule objectives and Kpoly(H ) :=
U D n
inf {− log w(a)} is the resource-bounded Kolmogorov complexity of the reusable
a∈S(Hn)
structure.
Assumption 23 (Sufficient-statistic condition) The history H is a sufficient statistic
n
for the strategy set S(H ) relative to the side information D :
n n
P (D | H , S(H )) = P (D | H ).
n n n n n
This assumption is natural in the packed-family setting, where H determines the posterior
n
allocation over modes and thereby determines the distribution of strategies.
Theorem 24 (Upper envelope at expectation level) Under Assumptions A1–A7 (As-
sumption 22) and the sufficient-statistic condition (Assumption 23):
E (cid:2) Sp(D ) (cid:3) ≤ I(H ; D ) + log (cid:0) 1/w(S(H )) (cid:1) ,
Dn n n n n
where Sp(D ) := log(T /T ) is the log-speedup from conditioning on D , I(H ; D ) is
n U UDn n n n
the mutual information between the history and the side information, and w(S(H )) is the
n
prior weight assigned to the strategy set.
Proof Step 1 (Pointwise ratio bound). By Assumptions A1–A7, there exists a∗ ∈ S(H )
n
such that for each realization of D :
n
Sp(D ) ≤ ρ + log w(a∗ | D ) − log w(a∗).
n Dn n
Step 2 (Expectations under sufficiency). Take expectations over D . The sufficient-
n
statistic condition ensures that a∗ does not depend on D in a way that correlates with the
n
posterior weights beyond the dependence through H , so Jensen’s inequality applies cleanly:
n
E (cid:2) log w(a∗ | D ) (cid:3) ≤ E (cid:2) sup log w(a | D ) (cid:3) .
Dn n Dn n
a∈S(Hn)
17
\end{verbatim}

\clearpage
\section*{Extracted page 18}
\begin{verbatim}
Beneventano and Poggio
Step 3 (Information bound). By the Kraft inequality applied to both prior and posterior:
E [ρ ] + E (cid:2) sup log w(a | D ) (cid:3) − log w(a∗) ≤ I(H ; D ) + log(1/w(S(H ))).
Dn Dn Dn n n n n
a∈S(Hn)
Step 4 (Combine). Collecting the bounds:
E (cid:2) Sp(D ) (cid:3) ≤ I(H ; D ) + log(1/w(S(H ))).
Dn n n n n
Which results inherit the sufficiency assumption. The sufficient-statistic condition
is required for this upper envelope and inherited by the approximation guarantee (Proposi-
tion 25), which references the upper envelope as a ceiling. The RI Identity and the four-term
decomposition are algebraic identities within the packed family and do not depend on the
sufficiency assumption.
Route A2 fallback. If the sufficiency assumption is relaxed, a maximal-inequality route
gives a weaker bound with an additional log |S(H )| term arising from a union bound over
n
the strategy set.
6.4 Approximation guarantee outside the packed family
The packed-family assumptions are stylized. Real tasks have soft mode posteriors and
approximate log-linear time scaling. The following proposition quantifies the degradation.
Proposition 25 (Approximation guarantee) Suppose the packed-family assumptions
(Assumption 20) hold except for two relaxations:
1. (Soft modes.) The posterior πM is δ-approximately one-hot: there exists s∗(D) ∈ [M ]
D n
with πM(s∗(D)) ≥ 1 − δ for some δ ∈ [0, 1).
D
(M)
2. (Approximate log-linearity.) The certified time satisfies log T (s) = log κ(M)+log c +
q sim
log t (M, s) − log q(s) + r(M, s, q) where |r(M, s, q)| ≤ η uniformly.
0
Then the four-term identity (Theorem 21) holds up to a remainder:
(cid:12) (cid:20) (cid:21)(cid:12)
(cid:12) (cid:12) (cid:12) E (cid:104) log T ( π M0) (S) − log T ( q M D ) (S) (cid:105) − log κ κ ( ( M M 0 ) ) + φ(M 0 , M) + G(M) − ε(M) (cid:12) (cid:12) (cid:12) ≤ R(δ, η, M n ),
where
R(δ, η, M ) ≤ 2δ log M + 2η.
n n
Proof The two relaxations contribute independently.
Soft modes. When πM is δ-approximately one-hot, H(πM) ≤ h(δ) + δ log(M − 1) ≤
D D n
h(δ) + δ log M , where h(δ) = −δ log δ − (1−δ) log(1−δ). For M ≥ exp(h(δ)/δ), we have
n n
H(πM) ≤ 2δ log M . Both the baseline and deployed terms contribute, giving a routing
D n
correction of at most 2δ log M .
n
Approximate log-linearity. The residual r(M, s, q) enters additively into each log-time
term. Since the identity subtracts two such terms, |r(M, s, q ) − r(M , s, π)| ≤ 2η.
D 0
Combining: R(δ, η, M ) ≤ 2δ log M + 2η.
n n
18
\end{verbatim}

\clearpage
\section*{Extracted page 19}
\begin{verbatim}
A Mathematical Framework for Building Agent Systems
When the bound is tight. For M = 8 modes, δ = 0.1, η = 0.05: R ≤ 0.42 + 0.10 = 0.52
n
nats. Against a typical decomposition value ∆ ≈ 3 nats (a 20x speedup), the remainder is
about 17%—meaningful but not invalidating. The bound is tight when the effective mode
count is small and the time model is approximately multiplicative.
7 Consequences for Model Choice and System Design
The decomposition is the accounting identity. Five consequences follow: a crossover formula,
a specialization-dominance checklist, verifier equalization, a multi-agent threshold, and a
Kelly-betting interpretation.
7.1 The crossover formula and effective mode count
Corollary 26 (Model-selection crossover) Under Assumption 20, model M achieves
1
lower expected certified log-time than model M whenever
2
κ(M )
2
log + φ(M , M ) − φ(M , M ) > G(M ) − G(M ) + ε(M ) − ε(M ) . (15)
0 1 0 2 2 1 1 2
κ(M )
1 (cid:124) (cid:123)(cid:122) (cid:125) (cid:124) (cid:123)(cid:122) (cid:125) (cid:124) (cid:123)(cid:122) (cid:125)
(cid:124) (cid:123)(cid:122) (cid:125) competence advantage information advantage routing penalty
cost savings
Proof Apply Theorem 21 to both models with the same baseline; subtract. The baseline
terms cancel.
Derivation of the exponential threshold formula. Equation (15) gives the general
crossover condition. To obtain a closed-form threshold in terms of the effective mode count
M , we make three additional assumptions:
eff
(S1) Verifier equalization: |G(M ) − G(M )| ≤ η for η small relative to log(κ /κ ) (Proposi-
1 2 2 1
tion 27), so the information terms approximately cancel.
(S2) Linear routing-mismatch scaling: the specialist’s routing mismatch grows approximately
linearly with the log of the effective mode count, ε(M ) − ε(M ) ≈ cˆ· log M , where
small big eff
cˆ > 0 is the per-mode routing penalty. This holds exactly when the specialist allocates
uniformly over M modes while the generalist routes according to the posterior, giving
eff
cˆ = 1; more generally, cˆ is an empirically estimated coefficient.
(S3) Net competence advantage φˆ independent of M : the within-mode competence difference
eff
φ(M , M ) − φ(M , M ) is approximately constant across task strata.
0 1 0 2
Under (S1)–(S3), the crossover condition becomes log(κ /κ ) + φˆ > cˆ · log M ,
big small eff
which gives the threshold:
(cid:18) (cid:19)
log(κ /κ ) + φˆ
M⋆ = exp big small . (16)
eff cˆ
For a cost ratio of 30x, comparable competence (φˆ ≈ 0), and binary verifier feedback
(cˆ = log 2, corresponding to one bit of routing penalty per mode doubling):
(cid:18) (cid:19)
log 30
M⋆ = exp = 2log 2 30 = 30.
eff log 2
19
\end{verbatim}

\clearpage
\section*{Extracted page 20}
\begin{verbatim}
Beneventano and Poggio
Tasks with fewer than 30 effective modes belong to the specialist. With richer verifier
feedback (cˆ > log 2), the threshold drops: a more informative verifier makes it harder for
the specialist to justify its cost advantage.
7.2 Specialization dominance
The crossover formula gives a threshold. The following checklist gives sufficient conditions
for specialization dominance that are individually checkable.
A smaller model dominates when:
(i) it is cheaper: κ < κ ;
small big
(ii) it is comparably competent on effective modes: φ(M , M ) ≥ φ(M , M );
0 small 0 big
(iii) its routing mismatch is no worse: ε ≤ ε ;
small big
(iv) verifier equalization holds: |G − G | is small relative to the cost ratio.
big small
When all four conditions hold, the specialist dominates on every channel simultaneously.
The checklist is a design diagnostic, not a discovery: the conditions individually assert the
small model is better on every dimension. The value lies in the decomposition into four
named, measurable channels and the explicit quantitative form of the crossover.
7.3 Verifier equalization
The information term G(M) = I (S : D ) is, in general, model-dependent: different models
π M
generate different probes and receive different verifier feedback. A key simplification occurs
when the verifier is rich enough to make G approximately model-independent.
Proposition 27 (Verifier equalization) Under Assumption 20, suppose both models’
residual posterior entropies are small:
E [H(πM1)] ≤ η , E [H(πM2)] ≤ η .
D D 1 D D 2
Then
|G(M ) − G(M )| ≤ η + η .
1 2 1 2
Proof G(M) = H(π)−E [H(πM)]. Since H(π) is model-independent, |G(M )−G(M )| =
D D 1 2
|E [H(πM2)] − E [H(πM1)]| ≤ η + η .
D D D D 1 2
Verifier equalization holds when the external checker provides enough feedback that both
models’ residual uncertainties about the latent mode are small. Formal proof assistants with
step-level error messages, comprehensive test suites with diagnostic output, and constraint
solvers with counterexample feedback are all settings where this is plausible. When verifier
equalization holds, the crossover formula simplifies: the information terms nearly cancel,
and the model-selection decision reduces to cost plus competence versus routing quality.
The verifier is the great equalizer: a rich external checker neutralizes the large model’s
informational advantage, leaving only per-step cost and competence to determine the winner.
20
\end{verbatim}

\clearpage
\section*{Extracted page 21}
\begin{verbatim}
A Mathematical Framework for Building Agent Systems
7.4 Multi-agent specialization threshold
The crossover formula compares two models. The multi-agent threshold extends the analysis
to teams of specialized agents versus a single generalist.
Assumption 28 (Multi-agent packed family with coordination cost) Consider a
packed latent-mode family with M modes and prior π. A single-agent design uses model M
n g
(generalist) with per-mode competence t (M , s), per-step cost κ(M ), and routing allocation
0 g g
q . A K-agent specialist-team design assigns each of K agents to a subset of modes. Agent
g
i has model M with per-mode competence t (M , s) and per-step cost κ(M ). The team
i 0 i i
allocation is q , and team(s) ∈ [K] denotes the agent assigned to mode s. The team
team
incurs a coordination cost c (K) ≥ 0 added to the expected log-time. The prior-weighted
coord
average per-step cost of the team is κ¯ :=
(cid:80)Mn
π(s) κ(M ).
team s=1 team(s)
Theorem 29 (Multi-agent specialization threshold) Under Assumption 28, suppose
G(M ) − G = γ for a known constant γ. Then the K-agent specialist team achieves
g team
strictly lower expected certified log-time than the single generalist if and only if
(cid:88) t 0 (M g , s) κ(M g )
π(s) log + log + (ε − ε ) > c (K) + γ.
g team coord
t (M , s) κ¯
s 0 team(s) team (cid:124) (cid:123)(cid:122) (cid:125)
(cid:124) (cid:123)(cid:122) (cid:125) pe (cid:124) r-ste (cid:123) p (cid:122) savin (cid:125) gs routing improvement
ϕspec: competence gain
Proof Apply Theorem 21 to both designs with the same prior-matched baseline.
For the generalist: ∆ = log(κ /κ(M )) + φ(M , M ) + G(M ) − ε . For the team,
g 0 g 0 g g g
the effective per-step cost is κ¯ and the effective competence is φ(M , M ) =
team 0 team
(cid:80)
π(s) log(t (M , s)/t (M , s)). The team incurs additive coordination cost:
s 0 0 0 team(s)
∆ = log(κ /κ¯ ) + φ(M , M ) + G − ε − c (K). Subtracting:
team 0 team 0 team team team coord
κ(M )
g
∆ − ∆ = log + ϕ − γ + (ε − ε ) − c (K),
team g spec g team coord
κ¯
team
(cid:80)
where ϕ = π(s) log(t (M , s)/t (M , s)) and the M terms cancel. The bicondi-
spec s 0 g 0 team(s) 0
tional follows because the four-term decomposition is an algebraic identity.
Corollary 30 (Multi-agent threshold under verifier equalization) If verifier equal-
ization holds, |G(M ) − G | ≤ η, then the specialist team dominates whenever
g team
κ(M )
g
ϕ + log + (ε − ε ) > c (K) + η.
spec g team coord
κ¯
team
This is a sufficient condition: the team may dominate even when the left-hand side falls
below c (K) + η, provided the actual information-gain difference γ is less than η.
coord
Proof From Theorem 29, the team dominates iff the LHS exceeds c (K) + γ. Since
coord
γ ≤ η by verifier equalization, the condition LHS > c (K) + η ≥ c (K) + γ is
coord coord
sufficient.
21
\end{verbatim}

\clearpage
\section*{Extracted page 22}
\begin{verbatim}
Beneventano and Poggio
Builder interpretation. The exact biconditional requires knowing γ exactly. In practice,
verifier equalization provides only |γ| ≤ η, and the corollary gives a conservative test: if
the combined savings from specialization exceed c (K) + η, the team is guaranteed to
coord
dominate. The gap between the conservative threshold and the exact threshold is at most 2η.
This split follows Radner’s [31] tradition of distinguishing exact equilibrium characterizations
from sufficient conditions for coordination gains.
7.5 The Kelly connection
The inner routing problem of the packed family—choosing the share allocation q to minimize
expected certified log-time given feedback D—is the same log-optimal portfolio problem
that appears in Kelly betting [23], universal portfolio theory [11], and algorithm-portfolio
scheduling [17]. The optimal allocation q∗(· | D) = π (·) maximizes expected log-speedup
D
just as the Kelly bet maximizes expected log-wealth growth. The routing mismatch ε =
E [KL(π ∥q )] is the expected log-wealth loss from suboptimal betting. The mathematical
D D D
identity is known; the contribution is its instantiation in the LLM agent setting, where it
prescribes posterior-matched routing as the compute-optimal strategy and decomposes the
gap from optimality into four measurable terms.
8 Application: Building Math Research Agents
The preceding sections develop the framework in domain-general terms. We now instantiate
it for building agents for formal theorem proving in Lean 4. This is not a toy example.
Teams at several organizations are actively building such systems, facing precisely the design
choices the framework addresses: which model, what orchestration, when to fine-tune, and
how to diagnose failures.
The framework yields concrete, differentiated prescriptions that vary across three iden-
tifiable task regimes within the same domain. The same framework applied to different
regimes within formal theorem proving produces opposing recommendations—a feature, not
a bug, since it reflects genuine structural differences in what makes these tasks hard.
8.1 Three-regime taxonomy
Formal theorem proving in Lean 4 is not a single task. The strategy space, the role of library
knowledge, and the nature of the creative bottleneck vary dramatically across problem types.
We identify three regimes, distinguished by their effective mode count M and the resulting
eff
implications for the decomposition terms.
Routine tactic application (M ≈ 3–5). These are lemmas solvable by a small number
eff
of well-known tactic patterns: induction followed by simplification, ring arithmetic, omega
for linear integer goals. The mode space is narrow, the verifier provides rich step-level
feedback, and a specialist trained on Lean tactics has seen thousands of similar proofs. The
decomposition predicts: the cost term log(κ /κ ) dominates because φ > 0 and ε
big small small
is small. This is the specialization-friendly regime.
Library-heavy formalization (M ≈ 8–15). These lemmas require knowing which
eff
specific Mathlib lemma to invoke—a retrieval problem embedded in a proof-construction
22
\end{verbatim}

\clearpage
\section*{Extracted page 23}
\begin{verbatim}
A Mathematical Framework for Building Agent Systems
Table 1: Three-regime taxonomy for formal theorem proving in Lean 4. Columns marked
[E] contain empirically estimated values reflecting domain judgment; these would
differ for other task domains. Columns marked [F] are framework-derived: they
follow from the decomposition structure given the mode count. A practitioner
in a different domain must estimate their own mode counts via the protocol in
Section 9.
Regime M [E] Transfer from coding [E] Recommended Likely bot- Recommended
eff
model class [F] tleneck [F] intervention [F]
Routine tactic 3–5 Strong Small fine-tuned Cost (κ) Cheaper model,
application specialist more attempts
Library-heavy 8–15 Partial Medium model + Routing (ε) Better retrieval,
formalization retrieval augmen- mode-specific
tation routing
Creative theo- 20+ Weak Frontier general- Competence Better training
rem search ist (φ) or informa- data, broader
tion (G) pre-training
problem. The mode space is moderately wide: multiple viable proof strategies exist, and the
key differentiator is whether the model can locate the right library fact. The decomposition
predicts: the routing term ε is the primary bottleneck. A specialist may have low κ but high
ε if its training did not cover the relevant library corner. This is the contested regime where
the recommendation is conditional.
Creative theorem search (M ≈ 20+). These problems require multi-step combinations
eff
of tactics, novel lemma inventions, and proof-state abstractions not well represented in
any existing training corpus. The decomposition predicts: either φ or G is the bottleneck,
because even a well-routed model needs to discover strategies not in its training distribution.
The cost advantage of a specialist is eroded by the routing penalty ε.
8.2 Worked examples
Example 1 (Undergraduate Lean Agent — Specialist Dominates) Team A builds
a Lean 4 agent for undergraduate-level theorem proving. They consider DeepSeek-Prover-V1.5
(7B specialist, κ = $0.001/call) versus Claude Opus 4 (frontier generalist, κ = $0.03/call)
A B
under a monthly budget of $500.
Decomposition estimation. From a 50-problem development set (satisfying the screening
heuristic: κ ratio = 30×, narrow domain, > 50 problems):
Model logκ0 φ G −ε ∆
κ
DeepSeek-7B +1.70 +0.90 +0.80 −0.10 +3.30
Routine (M ≈ 3)
eff Claude Opus −1.70 0.00 +0.90 −0.55 −1.35
Illustrative values under model parameters from published benchmarks and structural
predictions of the framework.
23
\end{verbatim}

\clearpage
\section*{Extracted page 24}
\begin{verbatim}
Beneventano and Poggio
Verdict: deploy the specialist. The specialist advantage is ∆ − ∆ = 4.65 nats,
A B
corresponding to exp(4.65) ≈ 105× speedup under equalized cost. All four terms favor
the specialist. The decomposition shows why: cost (+3.40 nats from 30× cheaper price),
competence (+0.90 nats from faster convergence on induction/simp patterns), and routing
(ε = 0.10 ≪ ε = 0.55). Information is approximately equal. Even if the competence
A B
advantage φ were zero, the cost advantage alone exceeds any plausible combination of routing
penalty and information gap in this narrow-mode regime.
What benchmark racing would say. “DeepSeek wins 88% of problems.” Correct but unin-
formative: it does not explain which lever drives the advantage, whether the advantage persists
under different cost assumptions, or what would need to change for the recommendation to
flip.
Example 2 (Research-Level Lean Agent — Conditional Verdict) Team B builds a
Lean 4 agent for research-level formalization: library-heavy proofs requiring knowledge of
specific Mathlib4 lemmas, with 8–15 effective modes.
Decomposition estimation. From a 100-problem development set:
Model logκ0 φ G −ε ∆
κ
DeepSeek-7B +1.70 +0.10 +1.60 −0.85 +2.55
Library (M ≈ 12)
eff Claude Opus −1.70 0.00 +1.80 −0.35 −0.25
The specialist advantage is ∆ − ∆ = 2.80 nats. But this headline number conceals a
A B
fragile structure.
Sensitivity. The specialist’s routing mismatch ε = 0.85 nats is large: it wastes attempts
A
on tactic families that do not match the library-heavy mode structure. If the specialist’s
routing could be improved to ε = 0.45 (via retrieval augmentation), the advantage grows
A
to 3.20 nats—robust. But on the roughly 20% of library-heavy problems where only Claude
succeeds (because the specialist has never seen the required Mathlib lemma), no amount of
cost advantage helps. The expected advantage, weighting by coverage, may be near zero or
negative.
Verdict: conditional. The framework identifies this as a close call and provides specific
conditions under which the recommendation flips: (i) if the team can reduce ε via retrieval
A
augmentation, the specialist dominates; (ii) if frontier-only coverage matters, the generalist is
safer; (iii) if cˆ is uncertain, the team should run the full pilot. The framework’s prescription
is “invest in routing improvement before committing to the specialist”—qualitatively different
from benchmark racing’s “just try both.”
What benchmark racing would say. “DeepSeek wins 65% of problems overall.” The
decomposition reveals this conflates the 80% where both succeed (specialist wins most by cost)
and the 20% where only the generalist succeeds.
8.3 Coding-to-math transfer
The framework provides a nuanced answer to whether training on code corpora helps with
formal theorem proving.
24
\end{verbatim}

\clearpage
\section*{Extracted page 25}
\begin{verbatim}
A Mathematical Framework for Building Agent Systems
Narrow-mode regime (routine tactic application). The dominant modes—pattern
matching on goal structure, symbolic rewriting, search over a small tactic vocabulary—
overlap substantially with code-generation modes. Post-training equivalence predicts that
coding-task fine-tuning transfers well.
Wide-mode regime (creative theorem search). The dominant modes—lemma inven-
tion, proof-state abstraction, library navigation across deep dependency chains—are not
represented in code corpora. The transfer analysis predicts: source-task fine-tuning on code
is insufficient because the target has modes not present in the source. Broader mathematical
pre-training matters more.
The opposing conclusions. The same framework yields: code training helps for routine
tactic application and code training does not help for creative theorem search. This reflects
the structural difference in mode overlap between coding and proving across narrow versus
wide mode regimes. A benchmark that averages across regimes would report “partial
transfer,” masking the fact that transfer is strong where it does not matter much and weak
where it matters most.
8.4 Four failure modes
The four-term decomposition provides a diagnostic framework for understanding why a math
agent is slow.
1. Missing structure (low G). The agent’s proof traces show that knowing the problem
identity does not help predict which strategy works. The model lacks the representations
to extract information from the proof state. Intervention: better training data covering
the missing modes.
2. Bad routing (high ε). The agent attempts tactics in the wrong order or wastes
attempts on strategies that cannot work for the given goal structure. Intervention:
improve retrieval or add mode-specific routing.
3. Expensive nodes (high κ). Each attempt costs too much, consuming budget before
the model can explore enough strategies. Intervention: switch to a cheaper model or
reduce per-step cost via distillation.
4. Bad graph shape (critical-path bottleneck). The orchestration forces sequential
steps when parallel branching would be more efficient. Intervention: restructure the
workflow.
9 Measurement Protocol and Empirical Program
The decomposition is useful only to the extent that κ, φ, G, ε, and M can be estimated
eff
from data. A domain-agnostic recipe, concrete estimators, and a pilot study design follow.
9.1 Domain-agnostic mode-discovery recipe
Before any domain-specific measurement, the builder must identify the latent modes. The
recipe has four steps, of which only Step 2 requires domain knowledge.
1. Collect successful solution traces from the target task. Run candidate models on a
development set. Pool all successful traces.
25
\end{verbatim}

\clearpage
\section*{Extracted page 26}
\begin{verbatim}
Beneventano and Poggio
2. Define a trace-level similarity metric appropriate to the domain. This is the one domain-
specific step: theorem proving uses edit distance on tactic sequences; code generation uses
AST diff; text generation uses embedding cosine similarity; retrieval-augmented systems
use tool-call sequence alignment.
3. Cluster with stability selection. Apply hierarchical clustering at multiple thresholds.
Retain the mode count that is stable across bootstrap resamples. Report the stability
range.
4. Validate predictiveness. Check that mode labels predict differential model performance:
success rates must differ meaningfully across clusters for at least one model. If mode
labels do not predict which model succeeds, the discovered modes are not design-relevant.
Steps 1, 3, and 4 are domain-agnostic. Step 2 requires domain knowledge but is a
well-posed question.
9.2 Screening heuristic
Before investing in the full measurement protocol, practitioners should apply a cheap
screening test.
Condition Action Rationale
κ ratio > 50× and nar- Deploy specialist, no pi- Cost term (> 3.9 nats) dom-
row domain lot needed inates all plausible φ, G, ε
κ ratio < 3× Benchmark-race with Cost term (< 1.1 nats) is
small comparison set comparable to estimation
noise
Fewer than ∼ 50 verifi- Do not run pilot Mode labeling requires suffi-
able problems cient within-cluster samples
κ ratio 3–50×, moderate Run the full pilot Decomposition’s marginal
complexity, ≥ 50 prob- value over benchmark rac-
lems ing is concentrated here
This converts the measurement protocol from “always do this” to “do this when it is
worth doing”—consistent with a framework that includes cost in its own decomposition.
9.3 Estimator summary
All estimators require the mode-labeling step above. The following table summarizes.
26
\end{verbatim}

\clearpage
\section*{Extracted page 27}
\begin{verbatim}
A Mathematical Framework for Building Agent Systems
Quantity Observable / Estimator Main Error Source
κ(M) Per-step cost from API pric- Variance in tokens/call. Use
ing: κˆ = price/1K tokens × median.
tokens/call; or latency.
φ(M , M) Within-mode competence Small sample per mode.
0
(cid:80)
gap: φˆ = πˆ(s) · Bootstrap CIs.
s
log(tˆ (M , s)/tˆ (M, s)), where
0 0 0
tˆ is median solve time per
0
mode.
G(M) Gˆ = H(πˆ) − 1 (cid:80) H(πˆM ). Un- Entropy estimation bias.
N i Di
der verifier equalization: check Miller–Madow correction:
|Gˆ(M 1 ) − Gˆ(M 2 )| < 0.5 nats; if Hˆ MM = Hˆ + M 2 n N −1.
so, use the average.
ε(M) εˆ = 1 (cid:80) K(cid:99)L(πˆM ∥qˆM). Fre- KL plug-in estimator biased
N i Di i
quency proxy if model does not upward. Laplace smooth-
output mode probabilities. ing.
(cid:16) (cid:17)
M M = exp 1 (cid:80) H(πˆMref) , Mode-labeling uncertainty
eff eff N i Di
where M is the cheapest can- propagates directly.
ref
didate model.
9.4 Pilot study design
We specify a pilot for comparing DeepSeek-Prover-V1.5 (7B specialist) against Claude Opus 4
(frontier generalist) on 300 Mathlib4 lemmas.
Benchmark construction. Extract candidate lemmas from Mathlib4 (commit pinned
after DeepSeek-Prover training cutoff). Classify into three strata by top-level tactic diversity:
routine (≤ 2 tactic families), library-heavy (specific Mathlib lemma invocations), and creative
(≥ 3 tactic families, no single tactic > 50% of steps). Sample 100 per stratum; split 50/50
into development and test sets.
Budget equalization. Equalize total API cost, not number of attempts. At κ =
A
$0.001/call and κ = $0.03/call with $0.064/problem: DeepSeek gets 64 attempts per
B
problem, Claude gets ≈ 56.
Mode labeling. Pool successful proof traces from both models on the development set.
Extract top-level tactic sequences; compute pairwise normalized edit distance; hierarchical
agglomerative clustering with average linkage. Sweep distance threshold and retain the
M stable across bootstrap resamples. Fix labels before computing any model-dependent
eff
quantities.
Five testable predictions. The pilot tests five qualitative predictions: (P1) DeepSeek
has ∼ 30× lower per-call cost; (P2) DeepSeek has better competence on routine problems
(φ > 0); (P3) Claude has higher information gain on creative problems; (P4) DeepSeek
27
\end{verbatim}

\clearpage
\section*{Extracted page 28}
\begin{verbatim}
Beneventano and Poggio
has lower routing mismatch on routine problems; (P5) a crossover exists: specialist wins
low-M , generalist wins high-M .
eff eff
Power analysis. Under log-normal certified times, a two-sample t-test on log-speedup
differences with effect size d = 0.5, α = 0.05, and 80% power requires ∼ 64 problems per
group. The 150-problem test set (50 per stratum) provides adequate power for stratum-level
comparisons.
Honest assessment criterion. If prediction accuracy in the contested regime (library-
heavy) is near 50% (within [45%, 55%]), the decomposition adds nothing beyond a cost-based
rule of thumb, and the paper reports this as a negative result.
10 Discussion and Conclusion
10.1 Limitations
The framework rests on the packed latent-mode family. The assumption that certified
time scales as t (s)/q(s)—that doubling the compute share halves the expected time—is a
0
geometric-trials model. The framework does not model correlated strategies, where success
on one approach provides partial information for another. The mode-labeling step is a
clustering judgment, not a mathematical object: different clustering thresholds yield different
mode counts, and the decomposition values change accordingly.
The framework assumes an external verifier. Tasks without binary or continuous
verification—open-ended generation, creative writing, policy analysis—fall outside its scope.
The framework addresses this limitation by not pretending to cover these tasks.
The four-term decomposition is exact only in the packed-family regime. Outside it, the
approximation guarantee (Proposition 25) provides a bound with remainder O(δ log M + η).
n
For tasks with many soft modes or strongly nonlinear time scaling, the remainder may
dominate.
The measurement protocol requires a pilot study. The pilot is cheap (estimated $725–$775
for the Lean study) but not free. The screening heuristic (Section 9.2) mitigates this by
identifying when the pilot is worth running.
10.2 Open questions
1. Joint optimization. The framework solves the model-choice sub-problem of the master
design program (2). Full joint optimization over all six coordinates—model, routing,
tools, memory, context, parallelism—remains open.
2. Concavity of mixed-source objective. The mixed-source post-training objective
∆ (m) involves a mutual-information term (concave in α) and a mismatch term (not
α→B
generically convex). Whether the combined objective is concave under natural structural
conditions is unknown.
3. Dynamic mode discovery. The current framework assumes modes are fixed. An
adaptive version that discovers and refines modes during deployment—updating the
decomposition as the system learns—would extend the framework to non-stationary task
distributions.
28
\end{verbatim}

\clearpage
\section*{Extracted page 29}
\begin{verbatim}
A Mathematical Framework for Building Agent Systems
4. Estimation of I(S ; Z ). The transfer characterization (Proposition 14) involves the
B A
mutual information between a target’s latent modes and a source’s learned representation.
Direct estimation requires model internals or representation probing and remains an open
problem.
5. Multi-agent coordination cost. The multi-agent threshold (Theorem 29) treats
c (K) as exogenous. Deriving c from the structure of the coordination protocol
coord coord
is an open problem connected to Radner’s [31] team theory.
10.3 What the framework gives and does not give
The framework gives the builder four things. First, a decomposable objective: instead of a
single opaque comparison, the builder sees four channels, each independently improvable.
Second, a crossover formula: the specialist-versus-generalist boundary is computable, not
guessable. Third, a transfer characterization: the builder knows when source-task training
helps, when to retrain the router, and when broader pre-training is needed. Fourth, a
measurement protocol: every quantity in the decomposition is estimable from a structured
pilot.
The framework does not give the builder a solved optimization problem. It does not
choose the orchestration architecture, the tool library, the memory system, or the parallelism
degree. It does not replace engineering judgment. It replaces unstructured engineering
judgment with structured engineering judgment: the builder still decides, but now knows
which channel to invest in and how much each channel contributes.
The gap between the framework and a full agent-design compiler is large. The framework
occupies a middle ground between theory that is too abstract to be actionable and engineering
that is too ad hoc to be systematic. Whether that middle ground is useful is ultimately an
empirical question—one the pilot study is designed to answer.
Appendix A. Continuous-Reward Extension
The main body restricts to binary verification: a task instance either produces a verified
solution (V = 1) or does not. This appendix generalizes the framework to continuous reward
n
signals R : X × Y → [0, R ], where R (x, y) measures the quality of output y on input x.
n max n
The generalization is structurally parallel to the binary case; the algebraic decomposition
carries through with three specific modifications.
A.1 Setup
Definition 31 (Continuous-reward task family) A continuous-reward task family is a
tuple (X , Y , R , B , P ) where X is the input space, Y is the output space, R : X ×Y →
n n n n n n n n n n
[0, R ] is the reward function (measurable, bounded), B is the budget, and P is the task
max n n
distribution over X . The binary case is recovered by setting R = 1{V = 1} with R = 1.
n n n max
A.2 Three modifications
Modification 1: Success probability → expected reward per mode. The mode-
conditional log-time log t (M, s) = log(1/p (M)) generalizes to the mode-conditional log-
0 s
29
\end{verbatim}

\clearpage
\section*{Extracted page 30}
\begin{verbatim}
Beneventano and Poggio
reward-rate: log tR(M, s) := − log r (M), where r (M) := E[R (X, Y ) | S = s, model M].
0 s s n
When R is binary, r = p and the expressions coincide.
n s s
Modification 2: Certified time → certified reward-time. The certified time T M(s)
q
is replaced by T R,M(s) := κ(M)·c /(q(s)·r (M)). In log-space: log T R,M(s) = log κ(M)+
q sim s q
log c + log tR(M, s) + log(1/q(s))—the same additive structure as the binary case.
sim 0
Modification 3: Posterior via reward observations. The posterior is updated via
Bayes’ rule on reward observations: πR,M(s) ∝ π(s) · (cid:81) p(R | S = s, M). The informa-
D j∈D j
tion gain becomes GR(M) = I (S : DR ).
π M
A.3 Generalized decomposition theorem
Theorem 32 (Four-term decomposition under continuous rewards) Under the
packed-family assumption with continuous rewards (Modifications 1–3), the expected certified
log-reward-time decomposes as:
∆R := E(cid:2) log T R,M0(S) − log T R,M(S) (cid:3) = log κ(M 0 ) + φR(M , M) + GR(M) − εR(M),
π qD κ(M) 0
where φR(M , M) := E [log(r (M)/r (M ))] is the competence advantage, GR(M) := I (S :
0 π S S 0 π
DR ) is the information gain from reward observations, and εR(M) := E [KL(πR,M∥q )] is
M D D D
the routing mismatch.
Proof The proof is identical in structure to the binary case (Theorem 21). The additive
log-structure of certified reward-time ensures the cross-entropy identity applies. The tower
property E [E [log tR(M, S)]] = E [log tR(M, S)] holds by the same argument as in the
D πR,M 0 π 0
D
binary case. The four terms emerge identically.
A.4 Binary recovery
Setting R = 1{V = 1}: r (M) = p (M), tR = t , T R,M = T M, πR,M = πM, and ∆R = ∆.
n n s s 0 0 q q D D
All terms reduce exactly to their binary counterparts.
A.5 Connection to optimal stopping
The routing problem in the continuous-reward case generalizes to a stopping/multi-armed-
bandit structure. This connects to Weitzman’s [42] Pandora’s Box (each mode is a “box” with
opening cost proportional to κ(M)/q(s) and expected reward r (M); the routing mismatch
s
measures the gap between the builder’s allocation and the reservation-value ranking) and to
Zilberstein’s [47] anytime-algorithm framework (the reward function is the natural quality
measure and the four-term decomposition tells the builder why the system’s reward-rate is
what it is). The framework identifies the stopping/routing structure and decomposes the
performance gap into four channels. It does not prove a new optimal stopping theorem: the
inner routing problem is known, and the outer model-choice problem is the contribution of
the four-term decomposition.
30
\end{verbatim}

\clearpage
\section*{Extracted page 31}
\begin{verbatim}
A Mathematical Framework for Building Agent Systems
Appendix B. Proofs of Supplementary Results
B.1 Full proof of Proposition 17 (orchestration equivalence bounds)
Proof Step 0: Pinsker bound. By Pinsker’s inequality applied pointwise: δ := ∥πA −
D D
(cid:113) √
πB∥ ≤ 2 KL(πA∥πB). Taking expectations and applying Jensen (concavity of ·):
D 1 √ D D
δ¯ := E [δ ] ≤ 2α.
D D
Proof of (a). Recall G − G = E [H(πB)] − E [H(πA)]. By Audenaert’s entropy
A B D D D D
continuity applied pointwise: |H(πA) − H(πB)| ≤ δ log(M /δ ). Taking expectations:
D D D n D
|G − G | ≤ E [δ log(M /δ )]. Since f(t) = t log(M /t) is concave on (0, M ], Jensen
A B D D √n D √ n n
gives E [f(δ )] ≤ f(δ¯) ≤ 2α log(M / 2α).
D D n
Proof of (b). Use the KL-difference decomposition: KL(p∥q)−KL(r∥q) = [H(r)−H(p)]+
(cid:80) [p(s) − r(s)] log(1/q(s)). Apply with p = πA, r = πB, take expectations. Bound the first
s D √D
term by part (a) and the second by δ¯ · log(1/q ) ≤ 2α log(1/q ).
min min
Proof of (c). If q ≥ 1/M , then log(1/q ) ≤ log M . Parts (a) and (b) are both
√ min n min n
O( α log M ).
n
B.2 Full proof of the alignment-certified upper envelope (Theorem 24)
See Section 6, §6.3 for the expectation-level statement. The pointwise argument follows the
structure in the main text:
Sp(D ) ≤ I (H : D ) + ρ + κ .
n poly n n D U
Full details are provided in the earlier presentation.
Appendix C. Scaling Laws (Compressed)
The four-term decomposition induces scaling predictions for how ∆ changes with model
scale and data scale.
Pre-saturation regime. When G(M) ≪ H(π) (the verifier feedback does not yet resolve
all modes), increasing data or model quality increases G approximately linearly. In this
regime, improving the information channel is the highest-leverage intervention.
Post-saturation regime. When G(M) ≈ H(π) (the verifier has identified the dominant
mode), further improvements come only from reducing κ or ε. This is the regime where
specialization dominates: the system already knows what to do, and the question is how
cheaply and accurately it can execute.
Crossover threshold scaling. The crossover threshold M⋆ = exp((log(κ /κ ) +
eff big small
φˆ)/cˆ) grows exponentially with the cost ratio and competence advantage, and shrinks
with richer verifier feedback (cˆ larger). As frontier models become cheaper (compressing
κ /κ ), the specialization-viable region shrinks. This predicts that the value of task-
big small
specific fine-tuning decreases over time as inference costs fall—a testable prediction.
31
\end{verbatim}

\clearpage
\section*{Extracted page 32}
\begin{verbatim}
Beneventano and Poggio
Appendix D. Notation and Assumption Audit
D.1 Notation reference
Symbol Definition
T = (X , Y , V , B ) Budgeted transductive task at scale n
n n n n
T (n, δ) Certified time-to-solution at confidence 1 − δ
A
∆ Certified log-speedup: log T − log T
A0 A
S ∈ [M ] Latent mode variable
n
π, π Prior and posterior over modes
D
q, q Baseline and deployed routing allocations
D
κ(M) Per-step cost of model M
t (M, s) Within-mode certified scale for model M on mode s
0
φ(M , M) E [log(t (M , S)/t (M, S))]: competence advantage
0 π 0 0 0
G(M) I (S : D ): information generation quality
π M
ε(M) E [KL(πM∥qM)]: routing mismatch
D D D
M exp(H(π )): entropy-effective mode count
eff D
M⋆ Crossover threshold: specialist dominates below, generalist
eff
above
d = (M, π, Ω, W, K, p) Design tuple
α Orchestration-equivalence parameter: E [KL(πA∥πB)]
D D D
η Post-training efficiency: transfer ratio from family A to B
A→B
c (K) Coordination cost for K-agent team
coord
(cid:80)
ϕ π(s) log(t (M , s)/t (M , s)): specialization gain
spec s 0 g 0 team(s)
D.2 Assumption audit
The following table traces which assumptions each result requires.
32
\end{verbatim}

\clearpage
\section*{Extracted page 33}
\begin{verbatim}
A Mathematical Framework for Building Agent Systems
Result Assumptions Required
RI Identity (Thm. 19) Packed family (Asm. 6), prior-matched base-
line
Four-term decomposition (Thm. 21) Model-dependent packed family (Asm. 20)
Upper envelope (Thm. 24) A1–A7 (Asm. 22), sufficient-statistic condition
(Asm. 23)
Approximation guarantee (Prop. 25) Relaxed packed family (δ-soft modes, η-
approximate log-linearity)
Crossover formula (Cor. 26) Model-dependent packed family
Crossover threshold (16) Additionally (S1)–(S3): verifier equalization,
linear mismatch scaling, constant competence
Verifier equalization (Prop. 27) Model-dependent packed family, small residual
posterior entropy
Multi-agent threshold (Thm. 29) Multi-agent packed family (Asm. 28), known
γ
Multi-agent corollary (Cor. 30) Additionally: verifier equalization with bound
η
Orchestration equivalence (Prop. 17) α-orchestration equivalence (Def. 16), common
mode bank
Post-training equivalence (Thm. 13) Packed family, shared latent space
Transfer characterization (Prop. 14) Target admits packed-family representation
Continuous-reward extension (Thm. 32) Packed family with Modifications 1–3
References
[1] A. Achille and S. Soatto. AI agents as universal task solvers: It’s all about time.
Preprint, 2026.
[2] R. Agarwal, M. C. Machado, P. S. Castro, and M. G. Bellemare. Contrastive behavioral
similarity embeddings for generalization in reinforcement learning. In ICLR, 2021.
[3] K. M. R. Audenaert. A sharp continuity estimate for the von Neumann entropy. J.
Phys. A: Math. Theor., 40(28):8127, 2007.
[4] P.-L. Bacon, J. Harb, and D. Precup. The option-critic architecture. In AAAI, 2017.
[5] N. Baram et al. Behavioral mode discovery via trajectory clustering. Preprint, 2024.
[6] M. Besta et al. Graph of thoughts: Solving elaborate problems with large language
models. Preprint, 2023.
[7] J. R. Birge and F. Louveaux. Introduction to Stochastic Programming. Springer, 2nd
edition, 2011.
[8] S. Borgeaud et al. Improving language models by retrieving from trillions of tokens. In
ICML, 2022.
[9] M. Chen et al. Evaluating large language models trained on code. Preprint, 2021.
33
\end{verbatim}

\clearpage
\section*{Extracted page 34}
\begin{verbatim}
Beneventano and Poggio
[10] H. Chernoff. Sequential design of experiments. Ann. Math. Stat., 30(3):755–770, 1959.
[11] T. M. Cover and J. A. Thomas. Elements of Information Theory. Wiley, 2nd edition,
2006.
[12] I. Csisza´r and J. Ko¨rner. Information Theory: Coding Theorems for Discrete Memoryless
Systems. Cambridge University Press, 2nd edition, 2011.
[13] K. Ellis et al. DreamCoder: Learning to code by writing programs in your sleep. In
PLDI, 2021.
[14] W. Fedus, B. Zoph, and N. Shazeer. Switch Transformers: Scaling to trillion parameter
models with simple and efficient sparsity. JMLR, 23(120):1–39, 2022.
[15] N. Ferns, P. Panangaden, and D. Precup. Metrics for finite Markov decision processes.
In UAI, 2004.
[16] M. Feurer et al. Efficient and robust automated machine learning. In NeurIPS, 2015.
[17] C. P. Gomes and B. Selman. Algorithm portfolios. Artif. Intell., 126(1–2):43–62, 2001.
[18] G. Grand et al. LILO: Learning interpretable libraries by compressing and documenting
code. Preprint, 2023.
[19] E. A. Hansen and S. Zilberstein. Monitoring and control of anytime algorithms: A
dynamic programming approach. Artif. Intell., 126(1–2):139–157, 2001.
[20] B. A. Huberman, R. M. Lukose, and T. Hogg. An economics approach to hard
computational problems. Science, 275(5296):51–54, 1997.
[21] F. Hutter, H. Hoos, and K. Leyton-Brown. Sequential model-based optimization for
general algorithm configuration. In LION, 2011.
[22] M. I. Jordan and R. A. Jacobs. Hierarchical mixtures of experts and the EM algorithm.
Neural Comput., 6(2):181–214, 1994.
[23] J. L. Kelly. A new interpretation of information rate. Bell Syst. Tech. J., 35(4):917–926,
1956.
[24] L. A. Levin. Universal sequential search problems. Probl. Peredachi Inf., 9(3):115–116,
1973.
[25] L. A. Levin. Randomness conservation inequalities. Inform. Control, 61(1):15–37, 1984.
[26] P. Lewis et al. Retrieval-augmented generation for knowledge-intensive NLP tasks. In
NeurIPS, 2020.
[27] Y. Li et al. Competition-level code generation with AlphaCode. Science, 378(6624):1092–
1097, 2022.
[28] J. Marschak and R. Radner. Economic Theory of Teams. Yale University Press, 1972.
34
\end{verbatim}

\clearpage
\section*{Extracted page 35}
\begin{verbatim}
A Mathematical Framework for Building Agent Systems
[29] M. Li and P. Vita´nyi. An Introduction to Kolmogorov Complexity and Its Applications.
Springer, 3rd edition, 2008.
[30] M. S. Pinsker. Information and Information Stability of Random Variables and Processes.
Holden-Day, 1964. (Russian original 1960.)
[31] R. Radner. Team decision problems. Ann. Math. Stat., 33(3):857–881, 1962.
[32] B. Ravindran and A. G. Barto. SMDP homomorphisms: An algebraic approach to
abstraction in semi-Markov decision processes. In IJCAI, 2003.
[33] S. Russell and E. Wefald. Do the Right Thing: Studies in Limited Rationality. MIT
Press, 1991.
[34] J. Schmidhuber. The speed prior: A new simplicity measure yielding near-optimal
computable predictions. In COLT, 2002.
[35] N. Shazeer et al. Outrageously large neural networks: The sparsely-gated mixture-of-
experts layer. In ICLR, 2017.
[36] C. Snell et al. Scaling LLM test-time compute optimally can be more effective than
scaling model parameters. In ICLR, 2025.
[37] R. J. Solomonoff. A formal theory of inductive inference, Part I. Inform. Control,
7(1):1–22, 1964.
[38] R. S. Sutton, D. Precup, and S. Singh. Between MDPs and semi-MDPs: A framework
for temporal abstraction in reinforcement learning. Artif. Intell., 112(1–2):181–211,
1999.
[39] M. E. Taylor and P. Stone. Transfer learning for reinforcement learning domains: A
survey. JMLR, 10:1633–1685, 2009.
[40] A. Wald. Sequential Analysis. Wiley, 1947.
[41] X. Wang et al. Self-consistency improves chain of thought reasoning in language models.
In ICLR, 2023.
[42] M. L. Weitzman. Optimal search for the best alternative. Econometrica, 47(3):641–654,
1979.
[43] L. Xu, F. Hutter, H. H. Hoos, and K. Leyton-Brown. SATzilla: Portfolio-based algorithm
selection for SAT. JAIR, 32:565–606, 2008.
[44] S. Yao et al. Tree of thoughts: Deliberate problem solving with large language models.
In NeurIPS, 2023.
[45] J. Zhang et al. AFlow: Automating agentic workflow generation. Preprint, 2024.
[46] K. Zheng et al. MiniF2F: A cross-system benchmark for formal Olympiad-level mathe-
matics. In ICLR, 2022.
35
\end{verbatim}

\clearpage
\section*{Extracted page 36}
\begin{verbatim}
Beneventano and Poggio
[47] S. Zilberstein. Using anytime algorithms in intelligent systems. AI Magazine, 17(3):73–
83, 1996.
36
\end{verbatim}

\end{document}
